\section{Numerical Precision and Limitations}\label{sec:numerical-stability}

Classical diagonalization caused no numerical difficulties throughout. The Hamiltonians involved are small — the largest is $5\times5$ — and \texttt{numpy.linalg.eigh} reproduces analytical eigenvalues to within $\approx10^{-15}$ across all parameter ranges, including near avoided crossings where the energy gap narrows. No ill-conditioning was detected at any of the interaction strengths studied.

The VQE experiments are more varied. All energy evaluations use exact statevector simulation, so finite-shot noise plays no role in the optimization. The only exception is the Bell state measurement, which used $1000$ shots per run; the observed spread ($\sigma\approx0.018$) is consistent with binomial sampling noise at that shot count. For all other VQE runs the dominant error source is optimizer convergence.

For the two-level Hamiltonian the single-parameter $R_y$ ansatz is exactly expressive, and residual errors of $\approx10^{-11}$ across the full $\lambda$ range reflect only the BFGS convergence tolerance. The landscape has a single minimum and a cold start from $\theta=0$ is sufficient throughout.

The two-qubit case is more nuanced. The $R_y$-CNOT-$R_y$ ansatz is equally expressive — it can represent any real two-qubit state — and for $\lambda\lesssim0.4$ it achieves machine-precision errors ($10^{-14}$--$10^{-11}$). However, at the avoided crossing near $\lambda\approx0.4$ the optimizer fails to follow the ground state to the lower branch and instead tracks the first excited state for all $\lambda>0.4$, producing errors of $\approx10^{-1}$. This is not a local-minimum problem in the conventional sense; the optimizer converges correctly at each $\lambda$, but to the wrong eigenstate. The failure is a branch-tracking limitation of independent cold-start optimization: without an explicit strategy for crossing between branches, the optimizer remains on the continuously connected upper branch.

The Lipkin $J=2$ system is where precision genuinely degrades. Mapping four particles to qubits embeds the problem inside a $16\times16$ Hilbert space that also contains $J=0$ and $J=1$ sectors; since the hardware-efficient ansatz carries no symmetry constraint, the optimizer must locate the physical ground state within the full unconstrained space. This is manageable away from the phase transition, but near $V_c\approx0.333$ the ground state changes character rapidly and the cost-function gradient becomes small, causing L-BFGS-B to stall. The resulting error of $\approx10^{-2}$ — roughly six orders of magnitude worse than the $J=1$ case — is not a generic optimizer failure; it reflects the structural mismatch between the Hamiltonian's symmetry and the unconstrained circuit. Multiple random restarts provide only partial relief. A symmetry-adapted ansatz confined to the $J=2$ subspace would address this directly, and is the most obvious improvement to pursue.

All results reported here assume exact statevector simulation. On actual NISQ hardware, gate errors, decoherence, and readout noise would add to the optimization difficulty already present in the $J=2$ case. The precision figures throughout this report should therefore be read as an optimistic upper bound, not as predictions of device-level performance.
